\chapter{The Saffron Gate}

\section*{Told by}
Merchant Zara, a Fhara woman who wears bells on her ankles and never haggles at the first offer. She says patience is a type of geometry, and silence is a type of currency, and the best bargains are the ones where both parties walk away feeling slightly robbed. She has never lost a negotiation. She has also never won one. She says that is the same thing.

\section*{The Tale}

A merchant came to the Saffron Gate. It was high as three men, carved with lions whose eyes followed the traveler, and guarded by two men who did not blink. Not once. Not even when dust blew into their faces. Not even when the merchant waved his hand in front of their eyes. They were not breathing. The merchant chose not to notice this.

\textit{``Toll is one truth,''} said the gate. The voice came from the lions, or from the guards, or from the stone itself. The merchant could not tell. He decided it did not matter.

He offered, \textit{``I am honest.''}

The gate did not open. The lions blinked. The merchant felt something behind his ribs shift---a small weight, like a coin sliding off a ledge.

He offered, \textit{``My goods are pure.''}

The gate did not open. The guards breathed in, once, together, like a bellows drawing air. The merchant felt another weight shift.

He offered, \textit{``I have never cheated a customer.''}

The gate swung open. Not slowly, not ceremoniously. Quickly, the way a door opens when someone on the other side has been holding the handle, waiting. The merchant stepped through. His legs were unsteady. His mouth was dry.

Inside, a woman sat on a stool. She had no bells on her ankles. She had no goods to sell. She had only her hands, folded in her lap, and her eyes, which were the same color as the gate's stone.

\textit{``You have just told three truths,''} she said. \textit{``The first was a boast. The second was a claim. The third was a lie. The gate only opens for lies. Did you not read the inscription?''}

The merchant looked back at the gate. There was an inscription. He had not read it. He had been too busy thinking about profit margins.

\textit{``Then what is the toll?''}

\textit{``The knowledge that you are a liar,''} said the woman. \textit{``Not the knowledge that you have lied. You knew that. The knowledge that you are a liar---that lying is not something you do, but something you are. That is the toll. That is the only toll the gate accepts.''}

She stood. She walked past him, through the gate, and vanished. The gate did not close. The merchant stood in the threshold for a long time, his goods on his back, his ledger in his pocket, his three truths still hanging in the air behind him.

He walked through. He did his business. He returned home. He never told anyone what he had learned. But he stopped believing his own stories. Not all of them. Just the ones where he was the hero.

The gate, they say, is still open. It is waiting for the next traveler who believes their own lies. The lions are still watching. The guards are still not breathing.

\section*{The Witch's Closing}
\textit{``A lie that you believe is not a lie. It is a cage. The gate opens from the inside, but only when you stop pretending you have the key. The key was never yours. The key was the door.''}

\section*{What Lurks in the Shadow of the Tale}

\subsection*{The Truth Gate}

A magical gate, carved into a mountain pass or standing alone in a desert, that opens only when a character speaks a falsehood they genuinely believe---not a deliberate lie, not a tactical misdirection, but a self-deception so deep that the character has mistaken it for the shape of their own face. The gate does not judge. The gate does not punish. The gate simply opens, and on the other side is a small, quiet room where a woman sits on a stool, waiting to tell you what you already know but have not yet said aloud.

\paragraph{Tags / Mechanics:}
\begin{itemize}
\item [TRUTH] [THRESHOLD] [REVELATION]
\item A character must declare a truth about themselves that is actually false. The GM may ask the player to describe the belief---not its origins, not its justifications, just the belief itself, stripped clean. If the GM judges that the character genuinely believes it (not knows it to be false, not suspects it might be shaky, but \textit{believes}), the gate opens.
\item Once opened, the character gains +1 die to all Insight rolls regarding that specific subject for the remainder of the campaign. They have seen their own blind spot. They cannot un-see it. The blind spot does not close; it simply becomes visible, like a crack in a wall that has always been there, waiting to be noticed.
\item The gate will not open for obvious lies, boasts, or rhetorical flourishes. It will not open for lies the character knows are lies. It will open only for the lies the character has mistaken for oxygen.
\item A character may attempt to open the gate more than once, but the toll increases. The second truth costs a memory. The third costs a name. The fourth---there is no recorded fourth. The woman on the stool does not speak of it.
\end{itemize}

\paragraph{Adventure Hook:}
A diplomat has been sent to open a sealed mountain pass---the only route through before winter closes the high roads. The pass is barred by a Saffron Gate. The gate requires a self-deception. The diplomat believes she is the best negotiator in the land, beloved by all, incapable of failure. She is none of these things. She has never been any of these things. The party must find a way to help her see the truth before the snow falls---or find someone else who can lie to themselves more convincingly. But the gate knows. The gate always knows.

\section*{Colophon}
Saffron bazaar, afternoon. The sun was the color of the gate's stone. Zara counted coins into eight piles. The ninth she left empty. Not as a toll. Not as an offering. As a reminder.

\textit{``For the lie I believe,''} she said. \textit{``And for the one I have not yet told.''} The empty pile did not move. The coins did not shift. But for a moment---just a moment---the ninth pile seemed deeper than the others. As if something had been placed there. As if something had always been there, waiting to be noticed.